\section{Conclusion and Future Work}\label{sec:conclusion}
Historical template compromise can remain consequential after a reference is renewed. In the evaluated pipelines, known old operators and repeated templates support reconstruction that passes a fresh-key matcher, despite rejection of unchanged old records and concealment of the fresh operator. At ten exposures, calibrated cross-image ASR in the matched CFP comparison is \MbfSmallAsr{} for MobileFaceNet and \RnetSmallAsr{} for ResNet50 at 128 bits; at 512 bits, the ranges are \MbfLargeAsr{} and \RnetLargeAsr{}. The separate original LFW MobileFaceNet study reaches \LfwBaselineAsr{} at 128 bits.

The same experiments place these security findings alongside recognition and storage. Both encoders support 16-byte and 64-byte references in place of 2048-byte raw embeddings. Although ResNet50 improves recognition, its higher utility does not confer reconstruction resistance under the tested exposure model. Compact-reference selection must therefore consider recognition, payload, and compromise jointly.

Further evaluation should measure latency, peak memory, energy, and operator storage for the complete pipeline on representative edge hardware. Additional encoders, larger impostor samples, independent correspondence review, and repeated model fitting would help assess the stability of the findings. Comparisons with alternative protection mechanisms should align genuine utility and attacker access. Image-generation and camera-level experiments would then be needed to assess attack interfaces beyond the digital embedding input studied here.